\documentclass[12pt, a4paper]{article}
\usepackage[utf8]{inputenc}
\usepackage[spanish]{babel}
\usepackage{amsmath, amsthm, amssymb, amsfonts}
\usepackage{enumitem}
\usepackage{geometry}
\usepackage{fancyhdr}
\usepackage{graphicx}
\usepackage{hyperref}
\usepackage{mathrsfs}
\usepackage{parskip}
\usepackage{setspace}
\usepackage{xcolor}

\usepackage{tikz}
\usepackage{pgfplots}
\pgfplotsset{compat=1.18}

% Configuración de la página
\geometry{left=2.5cm, right=2.5cm, top=3cm, bottom=3cm}
\pagestyle{fancy}
\fancyhf{}
\rhead{\textit{Cálculo para física}}
\lhead{\thepage}
\renewcommand{\headrulewidth}{0.4pt}

\setlength{\headheight}{14.49998pt}

% Definición de colores
\definecolor{azulrioplatense}{RGB}{0, 102, 204}
\definecolor{grisclaro}{RGB}{240, 240, 240}

% Estilo para teoremas y definiciones
\newtheoremstyle{rioplatense}
    {12pt}   % Espacio arriba
    {12pt}   % Espacio abajo
    {\itshape} % Cuerpo del teorema
    {0pt}    % Sangría
    {\bfseries} % Cabeza del teorema
    {.}      % Puntuación después de la cabeza
    { }      % Espacio después de la cabeza
    {}       % Cabeza especificada por el usuario

\theoremstyle{rioplatense}
\newtheorem{definicion}{Definición}[section]
\newtheorem{teorema}{Teorema}[section]
\newtheorem{ejemplo}{Ejemplo}[section]

\usepackage{url}
\usepackage{ulem}
%\usepackage{color}

%\usepackage[spanish]{babel}
%\usepackage{amssymb, amsthm}
%\usepackage{enumitem}
%\usepackage{mathrsfs}
\usepackage{float}
\usepackage{booktabs} % Para tablas elegantes
\usepackage{parskip}  % Para espaciado entre párrafos

% --- Configuración de estilo ---
\setlength{\parindent}{0pt}
\setlength{\parskip}{1em}
\renewcommand{\baselinestretch}{1.2} % Interlineado

% --- Comandos personalizados ---
\newcommand{\R}{\mathbb{R}}
\newcommand{\C}{\mathbb{C}}
\newcommand{\N}{\mathbb{N}}
\newcommand{\dom}{\text{dom}}

\newcommand{\fc}[1]{{\color{blue}{#1}}}

\title{Cálculo para física -- versión 1.0} 
\author{Franz Chouly}
\date{\today}

\begin{document}

\maketitle

\tableofcontents

%\documentclass[12pt, a4paper]{article}
%\usepackage[utf8]{inputenc}
%\usepackage[spanish]{babel}
%\usepackage{amsmath, amssymb, amsthm}
%\usepackage{graphicx}
%\usepackage{enumitem}
%\usepackage{mathrsfs}
%\usepackage{float}
%\usepackage{hyperref}
%\usepackage{booktabs} % Para tablas elegantes
%\usepackage{parskip}  % Para espaciado entre párrafos

% --- Configuración de estilo ---
%\setlength{\parindent}{0pt}
%\setlength{\parskip}{1em}
%\renewcommand{\baselinestretch}{1.2} % Interlineado

% --- Comandos personalizados ---
%\newcommand{\R}{\mathbb{R}}
%\newcommand{\C}{\mathbb{C}}
%\newcommand{\N}{\mathbb{N}}
%\newcommand{\dom}{\text{dom}}

%\title{\textbf{Apuntes de Funciones}}
%\author{Franz C.}
%\date{\today}

%\begin{document}
%\maketitle

\newpage

\section{Funciones}

\subsection{Lo mínimo}

\subsubsection{Definición de función}
Una \textbf{función} es una \textit{regla} que asigna a cada elemento de un conjunto \( A \) (el \textbf{dominio}) \textit{exactamente un} elemento de un conjunto \( B \) (el \textbf{codominio}). Formalmente, si \( f: A \to B \), entonces para cada \( x \in A \), existe un único \( y \in B \) tal que \( y = f(x) \).

\begin{itemize}
    \item \textbf{Notación:} \( f: A \to B \), \( x \mapsto f(x) \).
    \item \textbf{Analogía:} Una "máquina" que transforma un número en otro (ej: \( f(x) = x^2 \)).
\end{itemize}

\subsubsection{Ejemplos básicos}
\begin{itemize}
    \item \textbf{Polinomios:} \( p(x) = a_nx^n + \dots + a_0 \).
    \item \textbf{Raíz cuadrada:} \( f(x) = \sqrt{x} \), con \( x\geq 0 %\dom(f) = [0, +\infty) 
    \).
\end{itemize}

\subsubsection{Dominio de definición}
El \textbf{dominio} de \( f \) es el conjunto de valores \( x \) para los cuales \( f(x) \) \textit{está definido}. Ejemplos:
\begin{itemize}
    \item \( \displaystyle f(x) = \frac{1}{x} \): \( \dom(f) = \R \setminus \{0\} \).
    \item \( f(x) = \sqrt{x} \): \( \dom(f) = [0, +\infty) \).
\end{itemize}

\subsubsection{Gráficos}
La \textbf{gráfica} de \( f \) es el conjunto \( \{(x, f(x)) \mid x \in \dom(f)\} \). Herramientas útiles:
\begin{itemize}
    \item \textbf{Test de la recta vertical:} Si una recta vertical corta la gráfica en más de un punto, \textit{no es función}.
\end{itemize}

\subsubsection{Funciones pares e impares}
\begin{itemize}
    \item \textbf{Par:} \( f(-x) = f(x) \) (ej: \( f(x) = x^2 \)).
    \item \textbf{Impar:} \( f(-x) = -f(x) \) (ej: \( f(x) = x^3 \)).
\item Es importante\footnote{\textbf{Weyl, H.} (1964).
\emph{Symétrie et mathématique moderne}.
Nouvelle bibliothèque scientifique.
Flammarion, Paris, 156~p.
ISBN : 978-2-08-210141-7.}.
\end{itemize}

\subsubsection{Operaciones con funciones}
Dadas \( f \) y \( g \), se definen:
\begin{itemize}
    \item \textbf{Suma:} \( (f + g)(x) = f(x) + g(x) \), con \( \dom(f + g) = \dom(f) \cap \dom(g) \).
    \item \textbf{Producto:} \( (f \cdot g)(x) = f(x) \cdot g(x) \), mismo dominio.
    \item \textbf{Composición:} \( (f \circ g)(x) = f(g(x)) \), con \[ \dom(f \circ g) = \{x \in \dom(g) \mid g(x) \in \dom(f)\}. \]
\end{itemize}

\subsubsection{Función inversa}
Si \( f \) es \textbf{biyectiva}, existe \( f^{-1} \) tal que:
\[
f^{-1}(f(x)) = x \quad \text{y} \quad f(f^{-1}(y)) = y.
\]
\textbf{Ejemplos:}
\begin{itemize}
    \item \( f(x) = x \) es su propia inversa.
    \item \( f(x) = x^2 \) \textit{no es biyectiva} en \( \R \), pero sí en \( [0, +\infty) \), donde \( f^{-1}(y) = \sqrt{y} \).
\end{itemize}

\subsubsection{Monotonía}
\begin{itemize}
    \item \textbf{Creciente:} \( x_1 < x_2 \implies f(x_1) \leq f(x_2) \).
    \item \textbf{Decreciente:} \( x_1 < x_2 \implies f(x_1) \geq f(x_2) \).
\end{itemize}
\textbf{Ejemplos:}
\begin{itemize}
    \item \( f(x) = x^3 \) es creciente en \( \R \).
    \item \( f(x) = \displaystyle \frac{1}{x} \) es decreciente en \( (0, +\infty) \).
\end{itemize}

\subsection{Polinomios}

\subsubsection{Definición y ejemplos}
Un \textbf{polinomio} es una función de la forma:
\[
p(x) = a_nx^n + a_{n-1}x^{n-1} + \dots + a_0,
\]
donde \( a_i \in \R \) (o \( \C \)) y \( n \in \N \).

\subsubsection{Origen histórico}
Los polinomios surgieron en la \textbf{Edad Media} para resolver problemas geométricos (ej: cálculo de áreas) y ecuaciones algebraicas. Su estudio formal se desarrolló con el álgebra simbólica (siglos XVI-XVII).

\subsubsection{Propiedades fundamentales}
\begin{itemize}
    \item \textbf{Dominio:} \( \dom(p) = \R \) (o \( \C \)).
    %\item \textbf{Álgebra:} Los polinomios forman un \textit{anillo conmutativo} con la suma y el producto.
    \item \textbf{Composición:} Si \( p \) y \( q \) son polinomios, \( p \circ q \) también lo es.
    \item El conjunto de los polinomios con coeficientes en un cuerpo $\mathbb{K}$ (como $\mathbb{R}$ o $\mathbb{C}$) \textbf{forma un álgebra sobre $\mathbb{K}$}. Esto significa que, además de ser un espacio vectorial (donde podemos sumar polinomios y multiplicarlos por escalares), los polinomios se pueden multiplicar entre sí, y el resultado sigue siendo un polinomio. Por ejemplo, si $P(x) = 2x^2 + 3x + 1$ y $Q(x) = x + 2$, su suma $P(x) + Q(x) = 2x^2 + 4x + 3$ y su producto $P(x) \cdot Q(x) = 2x^3 + 7x^2 + 7x + 2$ son también polinomios. Además, la multiplicación cumple con las propiedades de \textbf{asociatividad}, \textbf{conmutatividad} y \textbf{distributividad}, lo que confirma que los polinomios son un álgebra conmutativa y unitaria (con $1$ neutro para la multiplicación).

\end{itemize}

\subsubsection{No-polinomios}
\begin{itemize}
    \item \textbf{Cociente:} \( \displaystyle \frac{p(x)}{q(x)} \) \textit{no es polinomio} (ej: \( \displaystyle {1}/{x} \)).
    \item \textbf{Inversa:} La inversa de \( p(x) = x^2 \) es \( \sqrt{x} \), que \textit{no es polinomio}.
\end{itemize}

\subsubsection{Interés}
Los polinomios son clave en:
\begin{itemize}
    \item \textbf{Series de Taylor.} %Aproximan funciones suaves cerca de un punto.

Las series de Taylor permiten aproximar funciones usando polinomios construidos a partir de las derivadas de la función en un punto. La idea central es que, si una función \( f(x) \) es derivable en \( x = a \), se puede expresar como \[ f(x) \simeq f(a) + f'(a)(x-a) + \frac{f''(a)}{2!}(x-a)^2 + \cdots.\] Cuantos más términos se incluyan, mejor será la aproximación cerca de \( a \). Por ejemplo, para funciones comunes como \( e^x \) o \( \cos(x) \), los primeros términos ya dan resultados sorprendentemente precisos para valores pequeños de \( x \).

Un caso clásico es la función exponencial \( e^x \), cuya serie de Maclaurin (Taylor en \( a=0 \)) al orden 2 es \[ e^x \simeq 1 + x + \frac{x^2}{2} .\] Para el coseno, en cambio, los términos impares se anulan, y al orden 2 queda \[ \cos(x) \simeq 1 - \frac{x^2}{2}.\] Estas aproximaciones son útiles en física, ingeniería y cálculo numérico, donde a menudo basta con polinomios simples para modelar comportamientos locales con buena precisión.

    \item \textbf{Interpolación de Lagrange.} %Pasa por puntos dados.

%\section{Interpolación de Lagrange: ajustar puntos con polinomios}

La interpolación de Lagrange permite construir un polinomio \( P(x) \) de grado \( n \) que pasa exactamente por \( n+1 \) puntos dados. A diferencia de las series de Taylor, no requiere derivadas, solo los valores de la función en los puntos elegidos. La fórmula es \( P(x) = \sum_{i=0}^n y_i \cdot L_i(x) \), donde \( L_i(x) \) son los polinomios base de Lagrange, definidos para anularse en todos los puntos excepto en \( x_i \). Este método es útil para aproximar funciones cuando solo se conocen algunos de sus valores, como en datos experimentales o tablas.

Por ejemplo, si interpolamos \( \cos(x) \) en los puntos \( x_0 = 0 \), \( x_1 = \pi/4 \) y \( x_2 = \pi/2 \) (con \( y_0 = 1 \), \( y_1 = \sqrt{2}/2 \) y \( y_2 = 0 \)), el polinomio de Lagrange de grado 2 resultante es
%\[
%P(x) = 1 \cdot \frac{(x - \pi/4)(x - \pi/2)}{(0 - \pi/4)(0 - \pi/2)} + \frac{\sqrt{2}}{2} \cdot \frac{(x - 0)(x - \pi/2)}{(\pi/4 - 0)(\pi/4 - \pi/2)} + 0 \cdot \frac{(x - 0)(x - \pi/4)}{(\pi/2 - 0)(\pi/2 - \pi/4)}.
%\]
Simplificando, se obtiene 
\( P(x) \simeq 1 - 0.6366x + 0.2251x^2 \), que coincide con \( \cos(x) \) en los puntos dados.
    
    \item Los \textbf{polinomios de Hermite} \(H_n(x)\) son una familia de polinomios ortogonales con respecto al peso \(e^{-x^2}\) en el intervalo \((-\infty, \infty)\). Estos polinomios, definidos por la relación de recurrencia \(H_{n+1}(x) = 2xH_n(x) - 2nH_{n-1}(x)\) con \(H_0(x) = 1\) y \(H_1(x) = 2x\), aparecen de manera natural en la resolución de la \textbf{ecuación de Schrödinger} para el \textbf{oscilador armónico cuántico}. En este contexto, las soluciones de la ecuación de Schrödinger independiente del tiempo para un potencial cuadrático \(V(x) = \frac{1}{2}m\omega^2x^2\) están dadas por las funciones de onda \(\psi_n(x) = N_n H_n(\alpha x)e^{-\frac{\alpha^2x^2}{2}}\), donde \(\alpha = \sqrt{\frac{m\omega}{\hbar}}\) y \(N_n\) es una constante de normalización. Los polinomios de Hermite permiten así describir los \textbf{estados cuánticos estacionarios} del oscilador armónico, con energías cuantizadas \(E_n = \hbar\omega(n + \frac{1}{2})\).

\end{itemize}

\subsection{Funciones trigonométricas}

\subsubsection{Fórmulas básicas}

\begin{enumerate}
\item \textbf{Definiciones básicas:}  

\[
\sin \theta = \frac{\text{opuesto}}{\text{hipotenusa}}, \quad  
\cos \theta = \frac{\text{adyacente}}{\text{hipotenusa}}, \quad  
\tan \theta = \frac{\text{opuesto}}{\text{adyacente}} = \frac{\sin \theta}{\cos \theta} .
\]

\item \textbf{Identidad pitagórica:}
\[
\sin^2 \theta + \cos^2 \theta = 1.
\]

%\item \textbf{Relación entre tangente y secante:}
%\[
%1 + \tan^2 \theta = \sec^2 \theta
%\]

%\item \textbf{Relación entre cotangente y cosecante:}
%\[
%1 + \cot^2 \theta = \csc^2 \theta
%\]

\item \textbf{Ángulos complementarios:}
\[
\sin\left(\frac{\pi}{2} - \theta\right) = \cos \theta, \quad
\cos\left(\frac{\pi}{2} - \theta\right) = \sin \theta.
\]
\item \textbf{Periodicidad:}
todas las funciones trigonométricas son \textbf{periódicas}, particularmente:
\begin{itemize}
    \item \( \sin(x + 2\pi) = \sin(x) \).
    \item \( \cos(x + 2\pi) = \cos(x) \).
\end{itemize}


\end{enumerate}

\subsubsection{Desarrollo histórico}
Las funciones trigonométricas surgieron de forma implícita en la antigüedad para estudiar astronomía y geometría (ángulos). %Hairer y Wanner (\textit{Analysis by Its History}) señalan su formalización en el siglo XVIII con Euler. *** Here for version 1.1. and add more about Fourier series, etc. Sine function: 
La función $\sen$ empezó a aparecer de forma más explicita en los trabajos de %Indian 
Brahmagupta (alrededor de 630) y luego %. Europe 
de Regiomontanus (1464). %**** 

%\subsubsection{Periodicidad}
%Todas las funciones trigonométricas son \textbf{periódicas}:
%\begin{itemize}
 %   \item \( \sin(x + 2\pi) = \sin(x) \).
 %   \item \( \cos(x + 2\pi) = \cos(x) \).
%\end{itemize}

%\subsubsection{Fórmulas imprescindibles}
%\begin{itemize}
%    \item \( \sin^2(x) + \cos^2(x) = 1 \).
%    \item \( \sin(a \pm b) = \sin(a)\cos(b) \pm \cos(a)\sin(b) \).
%f\end{itemize}

\subsubsection{Relaciones básicas y consecuencias}

\paragraph{Teorema de Ptolomeo, 150, Regiomontanus, 1467}
\[
\sin(a + b) = \sin a \cos b + \cos a \sin b
\]
\[
\cos(a + b) = \cos a \cos b - \sin a \sin b
\]

\paragraph{Fórmulas de ángulo doble}
\[
\sin(2x) = 2 \sin x \cos x
\]
\[
\cos(2x) = \cos^2 x - \sin^2 x
\]

\paragraph{Fórmulas de De Moivre}
Partemos de lo anterior y de:
\[
\sin(nx + x) = \sin(nx) \cos x + \cos(nx) \sin x
\]
\[
\cos(nx + x) = \cos(nx) \cos x - \sin(nx) \sin x
\]
Simplificando los argumentos:
\[
\sin((n+1)x) = \sin(nx) \cos x + \cos(nx) \sin x
\]
\[
\cos((n+1)x) = \cos(nx) \cos x - \sin(nx) \sin x
\]
Y iteramos. Obtenemos lo siguiente, con la notación
%\[
%(\cos a + i \sin a)^{n+1} = \cos((n+1)a) + i \sin((n+1)a)
%\]
\[
\binom{n}{k} = \frac{n!}{k!(n - k)!}.
\]

\paragraph{Fórmulas para múltiples ángulos}

%\subsubsection*{3x}
\[
\sin(3x) = \binom{3}{1} \sin x \cos^2 x - \binom{3}{3} \sin^3 x
\]
\[
\cos(3x) = \binom{3}{0} \cos^3 x - \binom{3}{2} \cos x \sin^2 x
\]

%\subsubsection*{4x}
\[
\sin(4x) = \binom{4}{1} \sin x \cos^3 x - \binom{4}{3} \sin^3 x \cos x
\]
\[
\cos(4x) = \binom{4}{0} \cos^4 x - \binom{4}{2} \cos^2 x \sin^2 x + \binom{4}{4} \sin^4 x
\]

%\subsubsection*{5x}
\[
\sin(5x) = \binom{5}{1} \sin x \cos^4 x - \binom{5}{3} \sin^3 x \cos^2 x + \binom{5}{5} \sin^5 x
\]
\[
\cos(5x) = \binom{5}{0} \cos^5 x - \binom{5}{2} \cos^3 x \sin^2 x + \binom{5}{4} \cos x \sin^4 x
\]

%\subsubsection*{6x}
\[
\sin(6x) = \binom{6}{1} \sin x \cos^5 x - \binom{6}{3} \sin^3 x \cos^3 x + \binom{6}{5} \sin^5 x \cos x
\]
\[
\cos(6x) = \binom{6}{0} \cos^6 x - \binom{6}{2} \cos^4 x \sin^2 x + \binom{6}{4} \cos^2 x \sin^4 x - \binom{6}{6} \sin^6 x
\]

\paragraph{Fórmula general para \(\cos(nx)\)}
Obtenemos:
\[
\cos(nx) = \sum_{k=0}^{\lfloor n/2 \rfloor} (-1)^k \binom{n}{2k} \cos^{n-2k} x \sin^{2k} x
\]
Más informalmente:
\begin{align*}
\cos(nx) = {} &
1 \cdot \cos^n x \\
&- \frac{n(n-1)}{2} \cos^{n-2} x \sin^2 x \\
&+ \frac{n(n-1)(n-2)(n-3)}{24} \cos^{n-4} x \sin^4 x \\
&+ \cdots
\end{align*}
%
%\[
%\cos(nx) =
%1 \cdot \cos^n x
%- \frac{n(n-1)}{2} \cos^{n-2} x \sin^2 x
%+ \frac{n(n-1)(n-2)(n-3)}{24} \cos^{n-4} x \sin^4 x
%+ \cdots
%\]

\paragraph{Fórmula general para \(\sin(nx)\)}
\[
\sin(nx) = \sum_{k=0}^{\lfloor (n-1)/2 \rfloor} (-1)^k \binom{n}{2k+1} \cos^{n-2k-1} x \sin^{2k+1} x
\]
Más informalmente:
\begin{align*}
\sin(nx) = {} &
n \cos^{n-1} x \sin x \\
&- \frac{n(n-1)(n-2)}{6} \cos^{n-3} x \sin^3 x \\
&+ \frac{n(n-1)(n-2)(n-3)(n-4)}{120} \cos^{n-5} x \sin^5 x \\
&+ \cdots
\end{align*}
Debemos estas formulas a de Moivre (1730) (ver también Euler, 1748).

\paragraph{Series}
En la formula $\cos(nx)$, hacemos el cambio de variable $X=nx$:
\begin{align*}
\cos(X) = {} &
1 \cdot \cos^n (X/n) \\
&- \frac{1(1-1/n)}{2/n^2} \cos^{n-2}(X/n) \sin^2 (X/n) \\
&+ \frac{1(1-1/n)(1-2/n)(1-3/n)}{24/n^4} \cos^{n-4} (X/n) \sin^4 (X/n) \\
&+ \cdots
\end{align*}
Lo mismo: 
\begin{align*}
\sin(X) = {} &
n \cos^{n-1} \left(\frac{X}{n}\right) \sin \left(\frac{X}{n}\right) \\
&- \frac{1 \cdot (1-1/n)}{6/n^2} \cos^{n-3} \left(\frac{X}{n}\right) \sin^3 \left(\frac{X}{n}\right) \\
&+ \frac{1 \cdot (1-1/n) \cdot (1-2/n) \cdot (1-3/n) \cdot (1-4/n)}{120/n^4} \cos^{n-5} \left(\frac{X}{n}\right) \sin^5 \left(\frac{X}{n}\right) \\
&+ \cdots
\end{align*}
Ahora hacemos, formalmente, siguiendo Jac. Bernoulli:
\[
n\rightarrow +\infty, \quad
\cos(X/n) \simeq 1, \quad
\sen(X/n) \simeq X/n, \quad 
(1-k/n) \simeq 1.
\]
Deducimos las series:
\begin{align*}
\cos x = 
1 
- \frac{x^2}{2!}
+ \frac{x^4}{4!} 
- \frac{x^6}{6!} 
+ \frac{x^8}{8!} 
- \frac{x^{10}}{10!} 
+ \cdots
\end{align*}
y
\begin{align*}
\sin x = 
x 
- \frac{x^3}{3!} 
+ \frac{x^5}{5!} 
- \frac{x^7}{7!} 
+ \frac{x^9}{9!} 
- \frac{x^{11}}{11!} 
+ \cdots
\end{align*}

%Teorema (Ptolemy, 150, Regiomontanus, 1467)
%\[
%\sen(a + b ) = \sen( a) \cos(b) + 
%\cos 
%\]
%Tambien
%\[
%sen(2x) =
%cos(2x) = 
%\]
%De Moivre
%\[
%sen((n+1)a 
%cos 
%\]
%Esto implica
%\[
%cos (3x) 
%\]
%Y más generalmente 
%\[
%cos (nx) 
%\]

\subsubsection{Series de Fourier}
Una clase importante de funciones periódicas pueden expresarse como suma de senos y cosenos (Gasquet-Witomski, \textit{Fourier Analysis and Applications}).

\paragraph{¿Qué son las series de Fourier?}
Las \textbf{series de Fourier} son una herramienta fundamental en matemática y física que permite descomponer una función periódica en una suma (posiblemente infinita) de funciones trigonométricas básicas: senos y cosenos. Esta idea, desarrollada por \textbf{Joseph Fourier} a principios del siglo XIX, surgió del estudio de la propagación del calor, pero hoy se aplica en áreas tan diversas como el procesamiento de señales, la acústica y la compresión de imágenes.

\paragraph{Contexto histórico}
Fourier presentó su teoría en 1807, revolucionando la forma en que los científicos entendían las ondas y los fenómenos periódicos. Su trabajo, inicialmente controvertido, demostró que funciones discontinuas podían representarse mediante series trigonométricas, algo que desafiaba las ideas matemáticas de la época.

\paragraph{Ejemplo clásico}
Consideremos la función continua $f(x) = x$ definida en el intervalo $[-\pi, \pi]$. Su serie de Fourier está dada por:
$$
f(x) \simeq \sum_{n=1}^{\infty} \frac{2(-1)^{n+1}}{n} \sin(nx).
$$
Aquí, cada término $\frac{2(-1)^{n+1}}{n} \sin(nx)$ representa una \textit{armónica} que, al sumarse, aproxima la función original. Cuantos más términos se incluyan, mejor será la aproximación.
%
%\section{Sumas parciales}
A continuación, se muestran las sumas parciales de la serie de Fourier. % para \( f(x) = x \) con 1, 2 y 3 términos.

\medskip

\noindent
{Suma parcial con 1 término (\( n = 1 \))}:
\[
S_1(x) %= \frac{2(-1)^{2}}{1} \sin(x) 
= 2 \sin(x).
\]
{Suma parcial con 2 términos (\( n = 1, 2 \))}:
\[
S_2(x) %= 2 \sin(x) + \frac{2(-1)^{3}}{2} \sin(2x) 
= 2 \sin(x) - \sin(2x).
\]
{Suma parcial con 3 términos (\( n = 1, 2, 3 \))}:
\[
S_3(x) = %2 \sin(x) - \sin(2x) + \frac{2(-1)^{4}}{3} \sin(3x) 
= 2 \sin(x) - \sin(2x) + \frac{2}{3} \sin(3x).
\]
%
%\section{Representación gráfica}
A continuación, se muestran las gráficas de las sumas parciales \( S_1(x) \), \( S_2(x) \), y \( S_3(x) \) junto con la función original \( f(x) = x \).

\begin{center}
\begin{tikzpicture}
\begin{axis}[
    axis lines = middle,
    xlabel = \(x\),
    ylabel = \(f(x)\),
    xmin=-3.5, xmax=3.5,
    ymin=-3.5, ymax=3.5,
    legend style={at={(1.05,0.5)},anchor=west}, % Légende à droite
    grid=both,
    grid style={line width=.1pt, draw=gray!10},
    major grid style={line width=.2pt,draw=gray!50},
    samples=200,
    domain=-pi:pi,
    restrict y to domain=-3.5:3.5,
    xtick={-3.14159, 0, 3.14159},
    xticklabels={$-\pi$, 0, $\pi$}
]

\addplot[blue, thick, line width=1.5pt] {x};
\addlegendentry{\(f(x) = x\)}

\addplot[red, dashed, line width=1.5pt] {2*sin(deg(x))};
\addlegendentry{\(S_1(x) = 2 \sin(x)\)}

\addplot[green, dotted, line width=1.5pt] {2*sin(deg(x)) - sin(2*deg(x))};
\addlegendentry{\(S_2(x) = 2 \sin(x) - \sin(2x)\)}

\addplot[orange, loosely dashed, line width=1.5pt] {2*sin(deg(x)) - sin(2*deg(x)) + (2/3)*sin(3*deg(x))};
\addlegendentry{\(S_3(x) = 2 \sin(x) - \sin(2x) + \frac{2}{3} \sin(3x)\)}

\end{axis}
\end{tikzpicture}
\end{center}

\paragraph{¿Por qué importa?}
Las series de Fourier son la base del análisis de frecuencias, esencial en telecomunicaciones, música y hasta en el diseño de filtros digitales. ¡Son el ``ADN matemático'' de las señales periódicas.

\subsection{Exponencial y logaritmo}

\subsubsection{Desarrollo histórico}
El estudio de la exponencial y el logaritmo se remonta a los siglos XVI-XVII, con aportes de Napier, Briggs y Euler. Hairer y Wanner (\textit{Analysis by Its History}) destacan su rol en la modelización de crecimiento y decaimiento.

\subsubsection{Propiedades básicas}
\begin{itemize}
    \item \textbf{Exponencial:} \( f(x) = a^x \) (\( a > 0 \)).
    \begin{itemize}
        \item \( a^{x+y} = a^x a^y \).
        \item \( \lim_{x \to +\infty} a^x = +\infty \) si \( a > 1 \).
    \end{itemize}
    \item \textbf{Logaritmo:} \( g(x) = \log_a(x) \) (\( a > 0, a \neq 1 \)).
    \begin{itemize}
        \item \( \log_a(xy) = \log_a(x) + \log_a(y) \).
        \item \( \log_a(a^x) = x \).
    \end{itemize}
\end{itemize}

\subsubsection{M\'as sobre exponencial}

\paragraph{El Triángulo de Pascal y la Fórmula del Binomio}

\paragraph{Introducción}
El Triángulo de Pascal es una disposición triangular de números que tiene aplicaciones en combinatoria, álgebra y teoría de probabilidades. Cada número en el triángulo es la suma de los dos números directamente encima de él. Además, los coeficientes binomiales \(\binom{n}{k}\) se pueden obtener directamente del triángulo.

\paragraph{Construcción del Triángulo de Pascal}
\begin{center}
\begin{tabular}{cccccccc}
& & & 1 & & & & \\
& & 1 & & 1 & & & \\
& 1 & & 2 & & 1 & & \\
1 & & 3 & & 3 & & 1 & \\
& \vdots & & \vdots & & \vdots & & \\
\end{tabular}
\end{center}

\paragraph{Propiedades del Triángulo de Pascal}
\begin{itemize}
    \item La suma de los elementos de la fila \(n\) es \(2^n\).
    \item Cada fila comienza y termina con 1.
    \item Los coeficientes son simétricos: \(\binom{n}{k} = \binom{n}{n-k}\).
\end{itemize}



\paragraph{Relación entre el triángulo de Pascal y la fórmula del binomio de Newton}
El \textbf{triángulo de Pascal} es una representación gráfica de los coeficientes binomiales $\binom{n}{k}$. Cada fila $n$ del triángulo corresponde a los coeficientes del desarrollo de $(a + b)^n$ en el binomio de Newton:
\begin{equation*}
(a + b)^n = \sum_{k=0}^{n} \binom{n}{k} a^{n-k} b^k
\end{equation*}
{Ejemplo para $n = 4$}.
%
El desarrollo de $(a + b)^4$ es:
%
\begin{equation*}
(a + b)^4 = a^4 + 4a^3b + 6a^2b^2 + 4ab^3 + b^4
\end{equation*}
%
Los coeficientes $1, 4, 6, 4, 1$ son los números en la cuarta fila del triángulo de Pascal.
%
\begin{center}
\begin{tabular}{cccccccccc}
&& & & 1 & & \\
&& & 1 & & 1 & \\
&& 1 & & 2 & & 1 \\
&1 & & 3 & & 3 & & 1 \\
1 && 4 & & 6 & & 4 && 1
\end{tabular}
\end{center}

\paragraph{Fórmula de recurrencia de los coeficientes binomiales}

Los coeficientes binomiales satisfacen la siguiente relación de recurrencia:
\begin{equation*}
\binom{n}{k} = \binom{n-1}{k-1} + \binom{n-1}{k}
\end{equation*}
con las condiciones iniciales:

\begin{equation*}
\binom{n}{0} = \binom{n}{n} = 1
\end{equation*}

%\subsection*{Prueba de la fórmula de recurrencia}

\paragraph{Demostración combinatoria}

Consideremos un conjunto de $n$ elementos. El número de subconjuntos de tamaño $k$, $\binom{n}{k}$, puede dividirse en dos casos:
\begin{itemize}
    \item Subconjuntos que incluyen un elemento específico (digamos, $x$): hay $\binom{n-1}{k-1}$ de estos.
    \item Subconjuntos que no incluyen el elemento $x$: hay $\binom{n-1}{k}$ de estos.
\end{itemize}
Por lo tanto, la relación de recurrencia se cumple:
\begin{equation*}
\binom{n}{k} = \binom{n-1}{k-1} + \binom{n-1}{k}
\end{equation*}
\paragraph{Demostración algebraica}
Usando la definición de los coeficientes binomiales:
\begin{align*}
\binom{n-1}{k-1} + \binom{n-1}{k} &= \frac{(n-1)!}{(k-1)!(n-k)!} + \frac{(n-1)!}{k!(n-1-k)!} \\
&= \frac{(n-1)!}{(k-1)!(n-k)!} + \frac{(n-1)!}{k!(n-k-1)!} \\
&= \frac{(n-1)!}{k!(n-k)!} \left( k + (n-k) \right) \\
&= \frac{(n-1)!}{k!(n-k)!} \cdot n \\
&= \frac{n!}{k!(n-k)!} \\
&= \binom{n}{k}
\end{align*}

\paragraph{Fórmula del binomio de Newton}
La fórmula del binomio de Newton expresa la expansión de \((a + b)^n\) como una suma de términos que involucran coeficientes binomiales:
\[
(a + b)^n = \sum_{k=0}^n \binom{n}{k} a^{n-k} b^k
\]
donde \(\binom{n}{k}\) es el coeficiente binomial, definido como:
\[
\binom{n}{k} = \frac{n!}{k!(n-k)!}
\]
La prueba de esta formula es elemental.
%\section*{Fórmula del binomio de Newton}
%
%La fórmula del binomio de Newton establece que, para cualquier entero no negativo $n$ y cualquier número real $a$ y $b$,
%
%\begin{equation*}
%(a + b)^n = \sum_{k=0}^{n} \binom{n}{k} a^{n-k} b^k
%\end{equation*}
%
%\section*{Prueba por inducción matemática}

\textbf{Caso base ($n = 0$)}
Para $n = 0$,
\begin{equation*}
(a + b)^0 = 1
\end{equation*}
Por otro lado,
\begin{equation*}
\sum_{k=0}^{0} \binom{0}{k} a^{0-k} b^k = \binom{0}{0} a^0 b^0 = 1
\end{equation*}
Por lo tanto, la fórmula se cumple para $n = 0$.

\textbf{Hipótesis de inducción}
Supongamos que la fórmula se cumple para algún entero $n = m$, es decir,
\begin{equation*}
(a + b)^m = \sum_{k=0}^{m} \binom{m}{k} a^{m-k} b^k
\end{equation*}

\textbf{Paso inductivo ($n = m + 1$)}
Queremos demostrar que la fórmula se cumple para $n = m + 1$, es decir,
\begin{equation*}
(a + b)^{m+1} = \sum_{k=0}^{m+1} \binom{m+1}{k} a^{(m+1)-k} b^k
\end{equation*}
Consideramos $(a + b)^{m+1} = (a + b)(a + b)^m$. Usando la hipótesis de inducción,
\begin{align*}
(a + b)^{m+1} &= (a + b) \sum_{k=0}^{m} \binom{m}{k} a^{m-k} b^k \\
&= a \sum_{k=0}^{m} \binom{m}{k} a^{m-k} b^k + b \sum_{k=0}^{m} \binom{m}{k} a^{m-k} b^k
\end{align*}
Reindexamos la segunda suma:
\begin{align*}
(a + b)^{m+1} &= \sum_{k=0}^{m} \binom{m}{k} a^{m+1-k} b^k + \sum_{k=1}^{m+1} \binom{m}{k-1} a^{m+1-k} b^k
\end{align*}
Combinamos las sumas:
\begin{align*}
(a + b)^{m+1} &= \binom{m}{0} a^{m+1} b^0 + \sum_{k=1}^{m} \left[ \binom{m}{k} + \binom{m}{k-1} \right] a^{m+1-k} b^k + \binom{m}{m} a^0 b^{m+1}
\end{align*}
Usamos la propiedad de los coeficientes binomiales: $\binom{m}{k} + \binom{m}{k-1} = \binom{m+1}{k}$,
\begin{align*}
(a + b)^{m+1} &= \binom{m+1}{0} a^{m+1} b^0 + \sum_{k=1}^{m} \binom{m+1}{k} a^{m+1-k} b^k + \binom{m+1}{m+1} a^0 b^{m+1}
\end{align*}
Lo que es equivalente a
\begin{equation*}
(a + b)^{m+1} = \sum_{k=0}^{m+1} \binom{m+1}{k} a^{(m+1)-k} b^k
\end{equation*}
Por el principio de inducción matemática, la fórmula del binomio de Newton se cumple para todo entero no negativo $n$.
%\subsection*{Ejemplo de Aplicación}
Por ejemplo, para \((a + b)^4\):
\[
(a + b)^4 = \binom{4}{0}a^4b^0 + \binom{4}{1}a^3b^1 + \binom{4}{2}a^2b^2 + \binom{4}{3}a^1b^3 + \binom{4}{4}a^0b^4
\]
\[
= 1 \cdot a^4 + 4 \cdot a^3b + 6 \cdot a^2b^2 + 4 \cdot ab^3 + 1 \cdot b^4
\]


\paragraph{Desarrollo de \((1 + \frac{x}{N})^N\) usando el binomio de Newton}

\paragraph{Fórmula compacta}
\[
\left(1 + \frac{x}{N}\right)^N = \sum_{k=0}^N \binom{N}{k} \left(\frac{x}{N}\right)^k
\]

\paragraph{Desarrollo de los 4 primeros términos}
\[
\left(1 + \frac{x}{N}\right)^N = \binom{N}{0} \left(\frac{x}{N}\right)^0 + \binom{N}{1} \left(\frac{x}{N}\right)^1 + \binom{N}{2} \left(\frac{x}{N}\right)^2 + \binom{N}{3} \left(\frac{x}{N}\right)^3 + \dots %\mathcal{O}\left(\frac{x^4}{N^3}\right)
\]
Sustituyendo los coeficientes binomiales:
\[
\left(1 + \frac{x}{N}\right)^N = 1 + N \cdot \frac{x}{N} + \frac{N(N-1)}{2} \cdot \frac{x^2}{N^2} + \frac{N(N-1)(N-2)}{6} \cdot \frac{x^3}{N^3} + \dots %\mathcal{O}\left(\frac{x^4}{N^3}\right)
\]
Lo que se simplifica a:
\[
\left(1 + \frac{x}{N}\right)^N = 1 + x + \frac{x^2}{2} \left(1 - \frac{1}{N}\right) + \frac{x^3}{6} \left(1 - \frac{1}{N}\right)\left(1 - \frac{2}{N}\right) + \dots %\mathcal{O}\left(\frac{x^4}{N^3}\right)
\]
\paragraph{Serie exponencial}
Tomando $N$ grande en la formula anterior (Euler):
\[
e^x =  1 + x + \frac{x^2}{2!} + \frac{x^3}{3!}  + \dots
\]
El método es peligroso. Por ejemplo así:
\[
1 = \frac{1}{N} + \ldots + \frac{1}{N} = \dots
\]
Pero la serie es valida. La prueba rigorosa se hara bien despues de Euler.

\paragraph{Origen del exponencial} Viene de un problema de geometria planteado por F. Debeaune (1601--1652), resuelto por Leibniz (1684). Esta vinculado igual a problemas prácticos (Euler, 1748), en demografía o economía, donde nos interesamos a expresiones como
\[
(1 + \omega )^N
\]
con $\omega$ pequeño y $N$ largo. Ver el libro de Hairer \& Wanner.
\subsection{Valor absoluto}

\subsubsection{Definición}
\[
|x| =
\begin{cases}
x & \text{si } x \geq 0, \\
-x & \text{si } x < 0.
\end{cases}
\]

\subsubsection{Propiedades}
\begin{itemize}
    \item \( |x| \geq 0 \) y \( |x| = 0 \iff x = 0 \).
    \item \( |xy| = |x||y| \).
    \item \( |x + y| \leq |x| + |y| \) (desigualdad triangular).
\end{itemize}

Para demostrar la \textbf{desigualdad triangular} en \(\mathbb{R}\), que establece que para cualesquiera \(a, b \in \mathbb{R}\) se cumple \(|a + b| \leq |a| + |b|\), procedemos de la siguiente manera:

Consideremos los casos posibles para los signos de \(a\) y \(b\):
\begin{itemize}
    \item Si \(a\) y \(b\) son ambos no negativos o ambos no positivos, entonces \(|a + b| = |a| + |b|\) directamente.
    \item Si \(a\) y \(b\) tienen signos opuestos, sin pérdida de generalidad, supongamos \(a \geq 0\) y \(b \leq 0\). Entonces, \(|a + b| \leq |a|\) o \(|a + b| \leq |b|\), dependiendo de cuál sea mayor. Entonces, tenemos \(|a + b| \leq \max\{|a|, |b|\} \leq |a| + |b|\), y la desigualdad se satisface.
\end{itemize}

Para una demostración formal, observamos que:
\[
-|a| \leq a \leq |a| \quad \text{y} \quad -|b| \leq b \leq |b|.
\]
Sumando estas desigualdades, obtenemos:
\[
-(|a| + |b|) \leq a + b \leq |a| + |b|.
\]
Por definición de valor absoluto, esto implica que \(|a + b| \leq |a| + |b|\), lo que completa la prueba.

\begin{center}
$\star$
Versión 1.2 -- 23 Marzo 26
$\star$
\end{center}

%\end{document}
\newpage

\section{Definición y concepto de límite}
\subsection{Introducción}
El concepto de \textbf{límite} es la piedra angular del análisis matemático. Permite estudiar el comportamiento de una función cerca de un punto, incluso cuando la función no está definida en ese punto. Intuitivamente, el límite describe el valor al que se acerca la función cuando su variable independiente se aproxima a un valor dado.

\subsection{Definición formal}
\begin{definicion}[Límite de una función]
    Sea $f: A \subseteq \mathbb{R} \to \mathbb{R}$ una función y $c \in \mathbb{R}$ un punto de acumulación de $A$. Decimos que $L \in \mathbb{R}$ es el \textbf{límite} de $f(x)$ cuando $x$ tiende a $c$, y lo denotamos como:
    $$
    \lim_{x \to c} f(x) = L,
    $$
    si para todo $\varepsilon > 0$, existe un $\delta > 0$ tal que para todo $x \in A$ con $0 < |x - c| < \delta$, se cumple que $|f(x) - L| < \varepsilon$.
\end{definicion}

\subsection{Unicidad del límite}
\begin{teorema}[Unicidad del límite]
    Si $\lim_{x \to c} f(x) = L_1$ y $\lim_{x \to c} f(x) = L_2$, entonces $L_1 = L_2$.
\end{teorema}
\begin{proof}
    Supongamos que $L_1 \neq L_2$. Tomemos $\varepsilon = \frac{|L_1 - L_2|}{2}$. Por definición de límite, existen $\delta_1$ y $\delta_2$ tales que:
    $$
    |f(x) - L_1| < \varepsilon \quad \text{y} \quad |f(x) - L_2| < \varepsilon,
    $$
    para $0 < |x - c| < \delta = \min\{\delta_1, \delta_2\}$. Sin embargo, por la desigualdad triangular:
    $$
    |L_1 - L_2| \leq |L_1 - f(x)| + |f(x) - L_2| < 2\varepsilon = |L_1 - L_2|,
    $$
    lo cual es una contradicción. Por lo tanto, $L_1 = L_2$.
\end{proof}

\subsection{Ejemplos}
\begin{ejemplo}
    Consideremos la función $f(x) = \frac{x^2 - 1}{x - 1}$. Observemos que $f$ no está definida en $x = 1$, pero podemos calcular su límite cuando $x \to 1$:
    $$
    \lim_{x \to 1} \frac{x^2 - 1}{x - 1} = \lim_{x \to 1} \frac{(x - 1)(x + 1)}{x - 1} = \lim_{x \to 1} (x + 1) = 2.
    $$
    Aquí, simplificamos la expresión algebraica para eliminar la indeterminación.
\end{ejemplo}

\section{Definición y concepto de continuidad}
\subsection{Definición formal}
\begin{definicion}[Continuidad en un punto]
    Sea $f: A \subseteq \mathbb{R} \to \mathbb{R}$ una función y $c \in A$. Decimos que $f$ es \textbf{continua en $c$} si:
    $$
    \lim_{x \to c} f(x) = f(c).
    $$
    Es decir, $f$ es continua en $c$ si el límite de $f(x)$ cuando $x \to c$ existe, es finito y coincide con $f(c)$.
\end{definicion}

\subsection{Teorema de Bolzano (valores intermedios)}
\begin{teorema}[Teorema de Bolzano]
    Sea $f: [a, b] \to \mathbb{R}$ una función continua en el intervalo cerrado $[a, b]$. Si $f(a)$ y $f(b)$ tienen signos opuestos (es decir, $f(a) \cdot f(b) < 0$), entonces existe al menos un $c \in (a, b)$ tal que $f(c) = 0$.
\end{teorema}
\begin{proof}
    La demostración se basa en el método de bisección. Consideremos el punto medio $m = \frac{a + b}{2}$. Si $f(m) = 0$, hemos terminado. Si no, $f(m)$ tiene el mismo signo que $f(a)$ o $f(b)$. Repetimos el proceso en el subintervalo correspondiente. Por el axioma de completez de los números reales, este proceso converge a un punto $c$ donde $f(c) = 0$.
\end{proof}

\subsection{Método de dicotomía para calcular un cero}
El teorema de Bolzano garantiza la existencia de un cero, pero no proporciona un método para encontrarlo. El \textbf{método de dicotomía} (o bisección) es un algoritmo iterativo que permite aproximar un cero de una función continua en un intervalo $[a, b]$ donde $f(a) \cdot f(b) < 0$.

\subsubsection{Algoritmo}
\begin{enumerate}
    \item \textbf{Inicialización:} Elige $a$ y $b$ tales que $f(a) \cdot f(b) < 0$.
    \item \textbf{Iteración:}
          \begin{itemize}
              \item Calcula el punto medio $m = \frac{a + b}{2}$.
              \item Si $f(m) = 0$, entonces $m$ es un cero de $f$.
              \item Si $f(a) \cdot f(m) < 0$, entonces el cero está en $[a, m]$. Actualiza $b = m$.
              \item Si $f(m) \cdot f(b) < 0$, entonces el cero está en $[m, b]$. Actualiza $a = m$.
          \end{itemize}
    \item \textbf{Repetición:} Repite el paso 2 hasta que $|b - a| < \varepsilon$, donde $\varepsilon$ es la precisión deseada.
\end{enumerate}

\subsubsection{Ejemplo}
Consideremos $f(x) = x^3 - x - 1$ en el intervalo $[1, 2]$. Claramente, $f(1) = -1$ y $f(2) = 5$, por lo que $f(1) \cdot f(2) < 0$. Aplicando el método de bisección:
\begin{itemize}
    \item Primera iteración: $m = 1.5$, $f(1.5) = 0.875$. El cero está en $[1, 1.5]$.
    \item Segunda iteración: $m = 1.25$, $f(1.25) = -0.234$. El cero está en $[1.25, 1.5]$.
    \item Continuamos hasta alcanzar la precisión deseada.
\end{itemize}

%\newpage

\subsection{Cálculo de límites clásicos}

\subsubsection{Indeterminaciones básicas}
Las indeterminaciones son formas que \textit{no permiten determinar} el límite directamente. Las más comunes son:
$$
\frac{0}{0}, \quad \frac{\infty}{\infty}, \quad 0 \cdot \infty, \quad \infty - \infty, \quad 0^0, \quad 1^\infty, \quad \infty^0.
$$

\subsubsection{Regla de L'Hôpital}
Si $ \lim_{x \to a} \frac{f(x)}{g(x)} $ es de la forma $ \frac{0}{0} $ o $ \frac{\infty}{\infty} $, y $ f $ y $ g $ son derivables cerca de $ a $, entonces:
$$
\lim_{x \to a} \frac{f(x)}{g(x)} = \lim_{x \to a} \frac{f'(x)}{g'(x)},
$$
siempre que el límite del lado derecho exista.

\textbf{Ejemplo:}
$$
\lim_{x \to 0} \frac{\sin(x)}{x} \quad \text{(forma } \frac{0}{0}\text{)}.
$$
Aplicando L'Hôpital:
$$
\lim_{x \to 0} \frac{\cos(x)}{1} = \cos(0) = 1.
$$

\subsubsection{Límites fundamentales}
\begin{enumerate}
    \item $\displaystyle \lim_{x \to 0} \frac{\sin(x)}{x} = 1 $.
    \item $\displaystyle \lim_{x \to 0} \frac{1 - \cos(x)}{x^2} = \frac{1}{2} $.
    \item $\displaystyle \lim_{x \to \infty} \left(1 + \frac{1}{x}\right)^x = e $.
    \item $\displaystyle \lim_{x \to 0} \frac{e^x - 1}{x} = 1 $.
\end{enumerate}

\textbf{Ejemplo de aplicación:}
$$
\lim_{x \to 0} \frac{\tan(x)}{x} = \lim_{x \to 0} \frac{\sin(x)}{x \cos(x)} = \frac{1}{1} = 1.
$$

\subsubsection{Trucos para resolver límites}

\paragraph{Racionalización}
Para límites con raíces, multiplicar por el conjugado:
$$
\lim_{x \to 0} \frac{\sqrt{x + 1} - 1}{x} = \lim_{x \to 0} \frac{(\sqrt{x + 1} - 1)(\sqrt{x + 1} + 1)}{x (\sqrt{x + 1} + 1)} = \lim_{x \to 0} \frac{x}{x (\sqrt{x + 1} + 1)} = \frac{1}{2}.
$$

\paragraph{Factorización}
Para polinomios, factorizar el numerador y denominador:
$$
\lim_{x \to 1} \frac{x^2 - 1}{x - 1} = \lim_{x \to 1} \frac{(x - 1)(x + 1)}{x - 1} = \lim_{x \to 1} (x + 1) = 2.
$$

\paragraph{Infinitésimos equivalentes}
Para $ x \to 0 $:
$$
\sin(x) \sim x, \quad \tan(x) \sim x, \quad 1 - \cos(x) \sim \frac{x^2}{2}, \quad e^x - 1 \sim x.
$$
\textbf{Ejemplo:}
$$
\lim_{x \to 0} \frac{\sin(3x)}{5x} = \lim_{x \to 0} \frac{3x}{5x} = \frac{3}{5}.
$$

\paragraph{Límites al infinito}
%
%\subsubsection{Polinomios}
El límite de un polinomio $ p(x) = a_nx^n + \dots + a_0 $ cuando $ x \to \infty $ depende del término dominante.

\begin{center}
$\star$
Versión 1.0 -- 12 Marzo 26
$\star$
\end{center}

\newpage

\section{Derivación}

Primero presentamos conceptos básicos, y luego unos teoremas importantes.

\subsection{Definición y concepto de derivada}
\subsubsection{Introducción}
La \textbf{derivada} de una función es uno de los conceptos fundamentales del análisis matemático. Representa la tasa de variación instantánea de una función en un punto y es la herramienta central del cálculo diferencial. Geométricamente, la derivada en un punto es la pendiente de la recta tangente a la curva en ese punto.

\subsubsection{Definición formal}
\begin{definicion}[Derivada de una función en un punto]
    Sea $f: A \subseteq \mathbb{R} \to \mathbb{R}$ una función y $c \in A$ un punto interior de $A$. Decimos que $f$ es \textbf{derivable en $c$} si existe el siguiente límite:
    $$
    f'(c) = \lim_{h \to 0} \frac{f(c + h) - f(c)}{h}.
    $$
    Este límite, si existe, se denomina \textbf{derivada de $f$ en $c$} y representa la pendiente de la recta tangente a la gráfica de $f$ en el punto $(c, f(c))$.
\end{definicion}

\subsubsection{Interpretación gráfica}
La derivada $f'(c)$ es la pendiente de la recta tangente a la curva $y = f(x)$ en el punto $(c, f(c))$. Si $f$ es derivable en $c$, la ecuación de la recta tangente es:
$$
y = f'(c)(x - c) + f(c).
$$

\begin{center}
\begin{tikzpicture}
    %\draw[->] (-2,0) -- (4,0) node[right]{$x$};
    %\draw[->] (0,-1) -- (0,4) node[above]{$y$};
    \draw[domain=0.5:3.5, smooth, variable=\x, azulrioplatense, thick] plot ({\x}, {0.5*(\x-2)^2 + 2});
    \draw[azulrioplatense, thick] (0.5,2) -- (3.5,2) node[right]{Recta tangente};
    \filldraw (2,2) circle (2pt) node[below right]{$(c, f(c))$};
    \node at (2.5,0.5) {$f'(c) = \text{pendiente}$};
\end{tikzpicture}
\end{center}

\subsubsection{Ejemplos de derivación de funciones elementales}
\begin{ejemplo}
    Calculamos las derivadas de algunas funciones elementales:
    \begin{itemize}
        \item Si $f(x) = k$ (constante), entonces $f'(x) = 0$.
        \item Si $f(x) = x^n$, entonces $f'(x) = n x^{n-1}$.
        \item Si $f(x) = \sin(x)$, entonces $f'(x) = \cos(x)$.
        \item Si $f(x) = e^x$, entonces $f'(x) = e^x$.
    \end{itemize}
\end{ejemplo}

\subsubsection{Reglas de derivación}
\begin{teorema}[Reglas de derivación]
    Sean $f$ y $g$ funciones derivables en $x$. Entonces:
    \begin{itemize}
        \item \textbf{Suma:} $(f + g)'(x) = f'(x) + g'(x)$.
        \item \textbf{Producto:} $(f \cdot g)'(x) = f'(x)g(x) + f(x)g'(x)$.
        \item \textbf{Cociente:} $\displaystyle \left(\frac{f}{g}\right)'(x) = \frac{f'(x)g(x) - f(x)g'(x)}{g(x)^2}$, si $g(x) \neq 0$.
        \item \textbf{Composición:} $(f \circ g)'(x) = f'(g(x)) \cdot g'(x)$.
    \end{itemize}
\end{teorema}

\begin{ejemplo}
    Calculemos la derivada de $f(x) = x^2 \sin(x)$:
    $$
    f'(x) = (x^2)' \sin(x) + x^2 (\sin(x))' = 2x \sin(x) + x^2 \cos(x).
    $$
\end{ejemplo}

\subsubsection{Monotonía y extremos}
\begin{definicion}[Función creciente/decreciente]
    Sea $f$ derivable en un intervalo $I$.
    \begin{itemize}
        \item $f$ es \textbf{creciente} en $I$ si $f'(x) \geq 0$ para todo $x \in I$.
        \item $f$ es \textbf{decreciente} en $I$ si $f'(x) \leq 0$ para todo $x \in I$.
    \end{itemize}
\end{definicion}

\begin{teorema}[Teorema de Fermat (puntos críticos)]
    Si $f$ tiene un extremo local en $c$ y es derivable en $c$, entonces $f'(c) = 0$.
\end{teorema}

\subsection{Teoremas de derivación}
\subsubsection{Teorema de Rolle}
\begin{teorema}[Teorema de Rolle]
    Sea $f: [a, b] \to \mathbb{R}$ una función continua en $[a, b]$ y derivable en $(a, b)$. Si $f(a) = f(b)$, entonces existe al menos un $c \in (a, b)$ tal que $f'(c) = 0$.
\end{teorema}

\subsubsection{Teorema de Lagrange (Incrementos Finitos)}
\begin{teorema}[Teorema del Valor Medio (Lagrange)]
    Sea $f: [a, b] \to \mathbb{R}$ una función continua en $[a, b]$ y derivable en $(a, b)$. Entonces existe al menos un $c \in (a, b)$ tal que:
    $$
    f'(c) = \frac{f(b) - f(a)}{b - a}.
    $$
\end{teorema}

\begin{proof}[Interpretación geométrica]
    El teorema afirma que existe un punto $c \in (a, b)$ donde la pendiente de la recta tangente a la curva $y = f(x)$ es igual a la pendiente de la recta secante que une los puntos $(a, f(a))$ y $(b, f(b))$.
\end{proof}

\begin{center}
$\star$
Versión 1.0 -- 16 Marzo 26
$\star$
\end{center}

\newpage

\section{Integración}
\subsection{Nociones básicas}
\subsubsection{Particiones}
\begin{definicion}[Partición de un intervalo]
    Sea $[a, b]$ un intervalo cerrado. Una \textbf{partición} $P$ de $[a, b]$ es un conjunto finito de puntos $x_0, x_1, \ldots, x_n$ tales que:
    $$
    a = x_0 < x_1 < \cdots < x_n = b.
    $$
    La \textbf{norma} de la partición, denotada $P$, es la longitud del subintervalo más grande:
    $$
    P = \max_{1 \leq i \leq n} x_i - x_{i-1}.
    $$
\end{definicion}

\subsubsection{Definición de integral de Riemann}
\begin{definicion}[Integral de Riemann]
    Sea $f: [a, b] \to \mathbb{R}$ una función acotada y $P = x_0, x_1, \ldots, x_n$ una partición de $[a, b]$. Para cada subintervalo $[x_{i-1}, x_i]$, definimos:
    $$
    m_i = \inf f(x) : x \in [x_{i-1}, x_i], \quad M_i = \sup f(x) : x \in [x_{i-1}, x_i].
    $$
    Las \textbf{sumas inferior y superior} de $f$ respecto a $P$ son:
    $$
    L(f, P) = \sum_{i=1}^n m_i (x_i - x_{i-1}), \quad U(f, P) = \sum_{i=1}^n M_i (x_i - x_{i-1}).
    $$
    Decimos que $f$ es \textbf{integrable según Riemann} en $[a, b]$ si:
    $$
    \sup_P L(f, P) = \inf_P U(f, P) = I,
    $$
    y en ese caso, $I$ se denomina la \textbf{integral de Riemann} de $f$ en $[a, b]$, denotada:
    $$
    \int_a^b f(x)  dx = I.
    $$
\end{definicion}

\subsection{Teorema fundamental del cálculo}
\subsubsection{Enunciado}
\begin{teorema}[Teorema Fundamental del Cálculo]
    Sea $f: [a, b] \to \mathbb{R}$ una función continua. Definimos $F: [a, b] \to \mathbb{R}$ como:
    $$
    F(x) = \int_a^x f(t)  dt.
    $$
    Entonces, $F$ es derivable en $(a, b)$ y se cumple:
    $$
    F'(x) = f(x) \quad \forall x \in (a, b).
    $$
    Además, si $G$ es cualquier primitiva de $f$ en $[a, b]$, entonces:
    $$
    \int_a^b f(x)  dx = G(b) - G(a).
    $$
\end{teorema}

\subsubsection{Aplicación al cálculo de primitivas}
\begin{ejemplo}
    Calculemos $\int_0^1 x^2  dx$ usando el Teorema Fundamental del Cálculo. Sabemos que $G(x) = \frac{x^3}{3}$ es una primitiva de $f(x) = x^2$. Por lo tanto:
    $$
    \int_0^1 x^2  dx = G(1) - G(0) = \frac{1^3}{3} - \frac{0^3}{3} = \frac{1}{3}.
    $$
\end{ejemplo}

\subsection{Métodos generales de integración}
\subsubsection{Cambio de variables}
\begin{teorema}[Cambio de variables]
    Sea $g: [\alpha, \beta] \to [a, b]$ una función derivable con derivada continua y biyectiva, y sea $f: [a, b] \to \mathbb{R}$ continua. Entonces:
    $$
    \int_a^b f(x)  dx = \int_{\alpha}^{\beta} f(g(u)) g'(u)  du.
    $$
\end{teorema}

\begin{ejemplo}
    Calculemos $\int_0^1 \sqrt{1 - x^2}  dx$ usando el cambio de variables $x = \sin(u)$. Entonces, $dx = \cos(u)  du$, y los límites de integración cambian a $u = 0$ y $u = \frac{\pi}{2}$:
    $$
    \int_0^1 \sqrt{1 - x^2}  dx = \int_0^{\frac{\pi}{2}} \sqrt{1 - \sin^2(u)} \cos(u)  du = \int_0^{\frac{\pi}{2}} \cos^2(u)  du.
    $$
    Usando la identidad $\cos^2(u) = \frac{1 + \cos(2u)}{2}$, obtenemos:
    $$
    \int_0^{\frac{\pi}{2}} \frac{1 + \cos(2u)}{2}  du = \frac{\pi}{4}.
    $$
\end{ejemplo}

\subsubsection{Integración por partes}
\begin{teorema}[Integración por partes]
    Sean $u$ y $v$ funciones derivables con derivadas continuas en $[a, b]$. Entonces:
    $$
    \int_a^b u(x) v'(x)  dx = u(x) v(x) \Big|_a^b - \int_a^b u'(x) v(x)  dx.
    $$
\end{teorema}

\begin{ejemplo}
    Calculemos $\int_0^1 x e^x  dx$ usando integración por partes. Tomamos $u(x) = x$ y $v'(x) = e^x$, de modo que $u'(x) = 1$ y $v(x) = e^x$. Aplicando la fórmula:
    $$
    \int_0^1 x e^x  dx = x e^x \Big|_0^1 - \int_0^1 e^x  dx = e - (e - 1) = 1.
    $$
\end{ejemplo}

\subsection{Métodos numéricos para aproximar una integral}
\subsubsection{Regla de los rectángulos}
Dada una partición $P = x_0, x_1, \ldots, x_n$ de $[a, b]$, la \textbf{regla de los rectángulos} aproxima la integral de $f$ como:
$$
\int_a^b f(x)  dx \approx \sum_{i=1}^n f\left(\frac{x_{i-1} + x_i}{2}\right) (x_i - x_{i-1}).
$$

\subsubsection{Regla del trapecio}
La \textbf{regla del trapecio} aproxima la integral de $f$ como:
$$
\int_a^b f(x)  dx \approx \sum_{i=1}^n \frac{f(x_{i-1}) + f(x_i)}{2} (x_i - x_{i-1}).
$$

\begin{ejemplo}
    Aproximemos $\int_0^1 x^2  dx$ usando la regla del trapecio con $n = 4$ subintervalos. La partición es $P = 0, 0.25, 0.5, 0.75, 1$. Aplicando la fórmula:
    $$
    \int_0^1 x^2  dx \approx \frac{0.25}{2} \left( f(0) + 2f(0.25) + 2f(0.5) + 2f(0.75) + f(1) \right) = 0.328125.
    $$
    El valor exacto es $\frac{1}{3} \approx 0.3333$, por lo que el error es pequeño.
\end{ejemplo}

\begin{center}
$\star$
Versión 1.0 -- 17 Marzo 26
$\star$
\end{center}

\newpage 

\section{Ecuaciones diferenciales}

\subsection{Definición de un problema de Cauchy}
\begin{definicion}[Problema de Cauchy]
    Un \textbf{problema de Cauchy} (o problema de valor inicial) para una ecuación diferencial ordinaria (EDO) de orden $n$ consiste en encontrar una función $y(t)$ que satisfaga:
    $$
    y^{(n)}(t) = f(t, y(t), y'(t), \ldots, y^{(n-1)}(t)),
    $$
    junto con las condiciones iniciales:
    $$
    y(t_0) = y_0, \quad y'(t_0) = y_1, \quad \ldots, \quad y^{(n-1)}(t_0) = y_{n-1},
    $$
    donde $f$ es una función dada, $t_0$ es un punto del dominio, y $y_0, y_1, \ldots, y_{n-1}$ son valores conocidos.
\end{definicion}

\begin{ejemplo}
    Un problema de Cauchy de primer orden tiene la forma:
    $$
    y'(t) = f(t, y(t)), \quad y(t_0) = y_0.
    $$
    Por ejemplo, el problema:
    $$
    y' = 2ty, \quad y(0) = 1,
    $$
    es un problema de Cauchy para una EDO lineal de primer orden.
\end{ejemplo}

\subsection{Ecuaciones lineales de primer orden}
\begin{definicion}[Ecuación lineal de primer orden]
    Una ecuación diferencial lineal de primer orden es de la forma:
    $$
    y'(t) + p(t)y(t) = g(t),
    $$
    donde $p(t)$ y $g(t)$ son funciones continuas en un intervalo $I$.
\end{definicion}

\subsubsection{Método de resolución}
La solución general de la ecuación $y' + p(t)y = g(t)$ se obtiene usando un \textbf{factor integrante} $\mu(t) = e^{\int p(t)  dt}$:
$$
y(t) = \frac{1}{\mu(t)} \left( \int \mu(t) g(t)  dt + C \right),
$$
donde $C$ es una constante arbitraria.

\begin{ejemplo}
    Resolvamos la ecuación $y' + 2ty = t$. El factor integrante es:
    $$
    \mu(t) = e^{\int 2t  dt} = e^{t^2}.
    $$
    Multiplicamos la ecuación por $\mu(t)$:
    $$
    e^{t^2} y' + 2t e^{t^2} y = t e^{t^2}.
    $$
    El lado izquierdo es la derivada de $e^{t^2} y$, por lo que:
    $$
    \frac{d}{dt} \left( e^{t^2} y \right) = t e^{t^2}.
    $$
    Integrando ambos lados:
    $$
    e^{t^2} y = \int t e^{t^2}  dt = \frac{1}{2} e^{t^2} + C.
    $$
    Despejando $y$:
    $$
    y(t) = \frac{1}{2} + C e^{-t^2}.
    $$
\end{ejemplo}

\subsection{Ejemplo del oscilador armónico}
\begin{definicion}[Oscilador armónico]
    El oscilador armónico simple está descrito por la EDO de segundo orden:
    $$
    m y''(t) + k y(t) = 0,
    $$
    donde $m$ es la masa, $k$ es la constante del resorte, y $y(t)$ es el desplazamiento en el tiempo $t$.
\end{definicion}

\subsubsection{Solución general}
La ecuación característica asociada es:
$$
m r^2 + k = 0 \implies r = \pm i \sqrt{\frac{k}{m}}.
$$
Por lo tanto, la solución general es:
$$
y(t) = A \cos\left(\sqrt{\frac{k}{m}} t\right) + B \sin\left(\sqrt{\frac{k}{m}} t\right),
$$
donde $A$ y $B$ son constantes determinadas por las condiciones iniciales.

\begin{ejemplo}
    Para $m = 1$, $k = 4$, y condiciones iniciales $y(0) = 1$, $y'(0) = 0$, la solución es:
    $$
    y(t) = \cos(2t).
    $$
\end{ejemplo}

\subsection{Ecuaciones no lineales en variables separables}
\begin{definicion}[Ecuación separable]
    Una ecuación diferencial de primer orden es \textbf{separable} si puede escribirse como:
    $$
    y' = f(t) g(y).
    $$
\end{definicion}

\subsubsection{Método de resolución}
Separamos variables e integramos:
$$
\int \frac{dy}{g(y)} = \int f(t)  dt.
$$

\begin{ejemplo}
    Resolvamos $y' = t y$. Separando variables:
    $$
    \int \frac{dy}{y} = \int t  dt \implies \ln|y| = \frac{t^2}{2} + C.
    $$
    Despejando $y$:
    $$
    y(t) = C e^{\frac{t^2}{2}}.
    $$
\end{ejemplo}

\subsection{Método de Euler explícito}
\begin{definicion}[Método de Euler]
    Dado el problema de Cauchy $y' = f(t, y)$, $y(t_0) = y_0$, el método de Euler aproxima la solución en los puntos $t_n = t_0 + n h$ ($n = 0, 1, \ldots$) mediante la iteración:
    $$
    y_{n+1} = y_n + h f(t_n, y_n),
    $$
    donde $h$ es el tamaño del paso.
\end{definicion}

\begin{ejemplo}
    Aproximemos la solución de $y' = t + y$, $y(0) = 1$ con $h = 0.1$ y $n = 2$:
    \begin{itemize}
        \item $t_0 = 0$, $y_0 = 1$.
        \item $t_1 = 0.1$, $y_1 = y_0 + h (t_0 + y_0) = 1 + 0.1 (0 + 1) = 1.1$.
        \item $t_2 = 0.2$, $y_2 = y_1 + h (t_1 + y_1) = 1.1 + 0.1 (0.1 + 1.1) = 1.22$.
    \end{itemize}
\end{ejemplo}

\begin{center}
$\star$
Versión 1.0 -- 18 Marzo 26
$\star$
\end{center}

\newpage

\section{EXTRA}

\newpage

\section{Sugerencias de lecturas}

La selección de obras que se presenta a continuación abarca desde reflexiones sobre la ciencia y las matemáticas hasta exploraciones filosóficas y neurocientíficas. Cada texto ha sido elegido por su capacidad para inspirar, desafiar y ampliar la comprensión de temas fundamentales en sus respectivos campos. Las obras están ordenadas alfabéticamente por autor para facilitar su consulta.

\vspace{1em}

\subsection*{Dahan-Dalmedico, A. \& Peiffer, J.}
Esta obra es una exploración fascinante del desarrollo histórico de las matemáticas, desde sus orígenes hasta su evolución en el contexto cultural, económico e institucional. Las autoras, Amy Dahan-Dalmedico y Jeanne Peiffer, ofrecen un análisis detallado de los temas clave que todo estudiante o docente enfrenta, como ecuaciones, espacio, límites y funciones. Lejos de presentar las matemáticas como un cuerpo estático de axiomas, este libro las muestra como una disciplina dinámica y en constante evolución.

\noindent
\textbf{Dahan-Dalmedico, A.} \& \textbf{Peiffer, J.} (1986).
\emph{Une histoire des mathématiques : Routes et dédales}.
Préface de Jean-Toussaint Desanti.
Points Sciences, Éditions du Seuil, Paris, 320~p.
ISBN : 978-2-02-009138-1.

\vspace{1em}

\subsection*{Hairer, E. \& Wanner, G.}
Este texto es una referencia esencial para quienes desean comprender el análisis matemático desde una perspectiva histórica. Ernst Hairer y Gerhard Wanner examinan cómo los conceptos fundamentales del análisis se han desarrollado a lo largo del tiempo, ofreciendo una visión enriquecedora para estudiantes y profesionales de las matemáticas.

\noindent
\textit{E. Hairer} and \textit{G. Wanner}, \emph{Analysis in historischer Entwicklung}. Übersetzt von Andreas Lochmann. Berlin: Springer (2011; Zbl 1247.26003)

\vspace{1em}

\subsection*{Lloyd, G.~E.~R.}
Geoffrey E. R. Lloyd, profesor emérito de filosofía y ciencia antiguas en la Universidad de Cambridge, ofrece en este libro una visión profunda de la ciencia griega, desde el siglo III a.C. hasta el siglo II a.C. La obra explora cómo surgió el pensamiento científico en la antigua Grecia y cómo este contexto histórico y cultural influyó en el desarrollo de la ciencia. Es una lectura imprescindible para historiadores de la ciencia y cualquier persona interesada en los orígenes del pensamiento científico.

\noindent
\textbf{Lloyd, G.~E.~R.} (1993).
\emph{Une histoire de la science grecque}.
Traduit de l'anglais par Charles Brunschwig.
Points Sciences, Éditions du Seuil, Paris, 448~p.
ISBN : 978-2-02-017765-8.

\vspace{1em}

\subsection*{Masson, S.}
Steve Masson, en esta obra accesible y práctica, explora cómo podemos optimizar el funcionamiento de nuestro cerebro. A través de ejemplos concretos y consejos basados en la neurociencia, el autor nos guía en el descubrimiento de estrategias para mejorar nuestras capacidades cognitivas. Este libro es ideal para quienes buscan entender mejor el cerebro y aplicar este conocimiento en su vida diaria.

\noindent
\textbf{Masson, S.} (2021).
\emph{Activer ses neurones}.
Odile Jacob, Paris, 256~p.
ISBN : 978-2-738-15150-6.

\vspace{1em}

\subsection*{Russell, B.}
Bertrand Russell, uno de los filósofos más influyentes del siglo XX, reúne en este libro una colección de ensayos que exploran el arte de filosofar. Con su estilo claro y penetrante, Russell aborda temas fundamentales de la filosofía, invitando al lector a reflexionar sobre el pensamiento crítico y la búsqueda de la verdad. Una lectura esencial para amantes de la filosofía y el pensamiento analítico.

\noindent
\textbf{Russell, B.} (1968).
\emph{The Art of Philosophizing, and Other Essays}.
Philosophical Library, New York, 119~p.

\vspace{1em}

\subsection*{Weyl, H.}
Hermann Weyl, matemático y físico teórico, presenta en este libro una reflexión sobre la simetría y su papel en las matemáticas modernas. La obra incluye un "álbum de imágenes" que ilustra la diversidad de formas y patrones simétricos en la naturaleza y el arte. Weyl combina rigor matemático con una sensibilidad estética, haciendo de este libro una lectura fascinante tanto para matemáticos como para amantes del arte y la ciencia.

\noindent
\textbf{Weyl, H.} (1964).
\emph{Symétrie et mathématique moderne}.
Nouvelle bibliothèque scientifique.
Flammarion, Paris, 156~p.
ISBN : 978-2-08-210141-7.

\end{document}

